\section{Experiments}
\subsection{Experimental Protocol}
The experiments are structured as a benchmark matrix with one detailed comparison case and several breadth checks. C-MAPSS FD001 is the most detailed case: it includes direct, fixed-depth, adaptive, generic-pretrained-LLM, and task-specific-LLM baselines on an official held-out test split. FD002--FD004 repeat the adaptive-versus-LLM comparison on additional official C-MAPSS splits. IMS repeats that comparison on three validation runs built from derived degradation stages. Paderborn extends the same adaptive-versus-LLM comparison into machinery state classification, while LBNL HVAC tests whether the looped model still transfers under stricter asset-level grouping.

Across all runs, the recurrent baselines use batch size 128, hidden size 128, and a maximum loop depth of six steps. The adaptive model uses the same backbone as the fixed-depth loop and differs only in the early-exit policy. The LLM comparison uses pretrained DistilBERT encoders over compact text summaries of the same telemetry windows, first as a generic pretrained model and then as a task-specific model with the top encoder block unfrozen for two epochs at a learning rate of $5\times10^{-5}$. The codebase also supports multitask supervision for maintenance action prediction and coarse subsystem attribution, along with stage-aware temporal validation for IMS.

\subsection{Metrics}
Classification results are reported with accuracy, macro F1, and weighted F1. Maintenance-specific reporting emphasizes critical-failure recall, action accuracy, subsystem accuracy, average loop depth, and measured latency in milliseconds per sample. Loop depth remains useful as an internal compute proxy, but the implementation now writes explicit latency benchmarks for validation and test artifacts as well.

\subsection{Reporting Strategy}
Because the datasets differ in supervision quality and split design, the paper reports them in three tiers. The first tier is the full FD001 comparison, used for the most detailed latency, depth, and error analysis. The second tier is the matched comparison matrix across FD002--FD004, IMS, and Paderborn, used to test whether the adaptive-versus-LLM ordering generalizes. The third tier is the HVAC transfer result, which is included to show domain reach without overstating an unfinished LLM comparison.

\subsection{Ablations}
The evaluation focuses on the core comparison between direct prediction, fixed-depth looping, and adaptive looping under a shared maximum depth. Useful follow-on ablations include the multitask loss weights, the stage-aware IMS label boundaries, maximum loop depth, and the exit threshold.

\subsection{Detailed Comparison Case}
FD001 is reported separately in the results section because it is the most complete test bed. It includes official test labels, the direct and fixed-depth recurrent baselines, and the clearest latency comparison against generic pretrained and task-specific LLM alternatives. That makes it the best place to inspect whether adaptive looping actually buys better maintenance behavior rather than just lower compute.
